% ============================================
% Response to Reviewer 1
% ============================================
\section*{Response to Reviewer 1}

We thank the reviewer for their positive assessment and for the remaining suggestions for improving clarity and reproducibility.

\subsection*{R1-1. Parental-community OD and pH}

\reviewercomment{The lower-OD parental community frequently wins in Nutr$+$, which may be consistent with strong acidification lowering endpoint OD while favoring Dominance. The reviewer asks whether this interpretation is worth mentioning in the main text and whether it can be supported by examining the relationship between parental-community OD and pH directly.}

We thank the reviewer for this suggestion, and we agree that the correlation is worth testing directly. We added the analysis to Supplementary Note~5 and expanded Supplementary Fig.~14 to show it. Endpoint OD$_{600}$ and pH are strongly positively correlated across parental-community replicates in Nutr$+$ ($n = 60$, Spearman $\rho = 0.82$, $p = 8.1 \times 10^{-16}$) and in Base ($n = 59$, $\rho = 0.79$, $p = 1.6 \times 10^{-13}$), but not in Nutr$-$ ($n = 60$, $\rho = 0.22$, $p = 0.091$). The lower-OD parental communities that win most often in Nutr$+$ are therefore also the more acidic ones, consistent with the reviewer's suggestion that strongly acidifying communities attain lower endpoint OD$_{600}$ yet dominate after mixing.

We have kept this in the Supplementary Information rather than adding it to the main Results. The association is correlational and does not establish that pH determines winner identity, and the main Results section already reports a directional test of the same idea: the acidic-versus-alkaline contrast between Base and Nutr$+$ (Extended Data Fig.~7). A second circumstantial line of support would lengthen that passage without strengthening the causal claim. The main Results section points readers to Supplementary Note~5 for these controls.

\subsection*{R1-2. Define Dominance in the Abstract}

\reviewercomment{Replace ``nutrient-dependent shift toward Dominance'' with wording that makes clear that Dominance means dominance of one parental community.}

We revised the Abstract to:

\mschange{``A similar nutrient-dependent shift toward dominance of one parental community also appears in taxonomically richer communities derived from natural environmental samples.''}

\subsection*{R1-3. Distinguish two failures of community cohesion}

\reviewercomment{Clarify that community cohesion can fail through two distinct routes: weak competitive exclusion between source communities and a lack of correlated species fates within parental communities.}

We revised the Introduction to clarify when species-level selection may occur during coalescence:

\mschange{``Together, these findings suggest that species-level selection may occur during coalescence when cross-community competitive exclusion is limited or when species-specific responses outweigh source-community identity in determining which species persist.''}

\subsection*{R1-4. Clarify the niche-partitioning sentence}

\reviewercomment{Revise the sentence to state more clearly that coexistence between species from different source communities via niche partitioning is more common than competitive exclusion.}

We revised the Introduction to:

\mschange{``Empirical work reports that coexistence between species from different source communities via niche partitioning is more common than competitive exclusion during large-scale biogeographic interchange in marine and terrestrial biomes.''}

\subsection*{R1-5. Define community composition vectors}

\reviewercomment{State explicitly that the community composition vectors contain relative abundances.}

We have added ``relative-abundance'' to the descriptions of the community composition vectors in the Results, Methods, and Supplementary Methods.

\subsection*{R1-6. Replace ``introduced'' with ``studied''}

\reviewercomment{Change ``introduced a generalized Lotka--Volterra model'' to ``studied'' or equivalent wording.}

We changed ``introduced'' to ``studied.''

\subsection*{R1-7. Clarify the 1,200 simulated coalescence events}

\reviewercomment{Clarify whether the four-community procedure is repeated 200 times to produce six pairwise coalescences per replicate and 1,200 total events, and explain why this structure is used rather than sampling 1,200 independent pairs directly.}

We clarified that the four-community procedure was repeated 200 times, with all six pairwise combinations tested within each replicate, yielding $200 \times 6 = 1{,}200$ total coalescence events. This structure mirrors the experimental workflow of assembling parental communities before testing pairwise coalescence.

\subsection*{R1-8. Remove repetitive ``Notably''}

\reviewercomment{The word ``Notably'' is used repetitively in the interaction-strength paragraph.}

We removed the repetitive uses of ``Notably.''

\subsection*{R1-9. Fig. 3 error-bar wording}

\reviewercomment{The Fig. 3 caption mentions error bars, but error bars are not visible in the figure.}

We removed the reference to error bars and clarified that the fractions are computed from 1,200 simulated coalescence events per value of $\mu$.

\subsection*{R1-10. Explain the expected nutrient-dependent shift}

\reviewercomment{Clarify the hypothesized connection between nutrient conditions, increased invasion resistance, stronger competition or environmental feedbacks, and the predicted shift in coalescence outcomes.}

We clarified the hypothesized link between nutrient supply, invasion resistance, competition, and environmental feedbacks and stated that stronger effective interactions are expected to shift coalescence outcomes from Mixture toward Dominance.

\subsection*{R1-11. Clarify categorical-classification wording}

\reviewercomment{Replace ``before applying category labels'' with ``before classifying outcomes categorically'' or equivalent wording.}

We changed the phrase to ``before classifying outcomes categorically.''

\subsection*{R1-12. Restore paragraph flow around alternative explanations}

\reviewercomment{Present the nutrient-gradient result and its consistency with the theoretical prediction first, then discuss in a separate paragraph whether changes in parental-community properties could provide alternative explanations.}

We reordered the section so that the nutrient-dependent results are followed immediately by their theoretical interpretation. A separate paragraph then addresses alternative explanations, including the existing biomass-heterogeneity and initial-richness controls.

\subsection*{R1-13. Explain gray and colored marks in Figs. 4C and 6B}

\reviewercomment{Clarify what the gray versus colored points and histogram bars represent in Figs. 4C and 6B.}

Because parental communities A and B have no intrinsic ordering, each event is shown twice: once in color with the original labels and once in gray with the labels reversed. We clarified this convention in the captions of Figs.~1, 4, 5, and 6.

\subsection*{R1-14. Clarify the dominant-species PDI control}

\reviewercomment{State explicitly that the predictive power of dominant-species interaction outcomes is not simply a consequence of including the dominant species in the PDI calculation.}

We revised the Results accordingly.

\subsection*{R1-15. Define or remove CLS in Fig. 5C}

\reviewercomment{The abbreviation ``CLS'' is used in Fig. 5C but is not defined in the figure or caption.}

We defined CLS as community-level selection in the Fig.~5 caption.

\subsection*{R1-code. Code availability}

\reviewercomment{The linked GitHub repository is not sufficient to conduct all analyses and does not include the code used to produce the figures. Additional analysis and figure-generation code should be made publicly available.}

We have expanded the public repository to include the full analysis pipeline and
figure-generation code for most main-text panels; a small number of experimental
schematics and standalone displays are not generated programmatically. We have
deposited an archived release at Zenodo
(DOI: \href{https://doi.org/10.5281/zenodo.21696637}{10.5281/zenodo.21696637})
and updated the Code Availability statement accordingly.
